\chapter{Goldfish Don't Bark}
'The nice thing about goldfish is that they don't bark!' said Granny, as she stood back from the goldfish tank to admire the pretty gold, orange, yellow and striped aquarium fish as they glided about their silent, glassed-in enclosure.
\par
'They don't make any noise at all,' said her granddaughter Koki, who had joined her in front of the tank. 'They don't bark like dogs, they don't bray like donkeys, they don't sing like birds. How do they talk to each other?'
\par
'They can communicate without talking' said Granny. 'Do you ever see them bump into each other? Humans do a lot of talking, and ever they keep bumping into each other.'
\par
Granny had acquired the aquarium recently. For over a month, she kept a dog, a scruffy-looking terrier, but the dog would bark all day and all night. In barked other dogs, it barked at squirrels, it barked at mice (but never chased them), it barked at shadows, it barked with doors creaked, it barked at autumn leaves falling from the chestnut trees.
\par
Granny couldn't stand so much barking. So, she gave the dog to her sister, who lived on the other side of the hill. Her sister liked dogs that barked. Some people do. Granny preferred goldfish.
\par
What did Granny look like? Well she looked like everyone's nice Granny-kind, while haired, a little plump, wearing an old pair of spectacles. And she smelled nice, like freshly baked bread.
\par
And what did Koki like?
\par
She looked like everyone's sweetest, prettiest little sister. What more can I say?
\par
She had two homes-one in the city, with her parents, and another here in foothills with her grandmother, whose like cottage was in a forest of Himalayan oak and maple trees. Koki spent as much time as she could with her grandmother, over the summer holidays, the winter break, even on Holi or Dussehra. In India, children gets lots of homework, but they also get a lot of holidays in which to do it!
Koki was doing her homework when she noticed a large ginger-coloured cat sitting at an open window, staring a gold fish at the other end of the room. On warm summer days, the window was usually kept wide open. It opened on to the trees, and sometimes a bird from the forest a-wagtail or a bulbul or a scarlet minivet-would hop in, explore the room and then return to the forest.
\par
Granny didn't mind these visitors-in fact, she welcomed them-and she didn't mind the cat, either, at least it didn't bark!
\par
But Granny had forgotten that cats like fish. Not because they are pretty, but because, for a cat, even a small fish is a tasty dish.
\par
Koki was noticed the cat. She stroked and petted it and the cat went away. Koki closed the window.
\par
But the next day, the cat was back, and this time it managed to conceal itself behind the window curtains. Koki had gone into the kitchen to help Granny make lunch. Suddenly, they heard a crash in the front room, and a strangled cry from the cat. Granny and Koki dropped what they were doing and ran to see what had happened.
\par
They found the cat upside down in the goldfish tank, half in and half out of the water. It had climbed up the stand, but had been unable to balance itself on the edge of the tank and fallen in. The tank was half empty and two goldfish lay on the carpet, flopping about helplessly.
\par
Koki got the cat out of the tank with some difficultly. It showed no gratitude and fled from the house, growling in shock and anger. It wouldn't be back for some time!
\par
But Granny had lost two goldfish, and now only three remained in the tank. So the next day, when they visited the Mussoorie Mall, they bought three new goldfish. One was a very pretty, deep pink little fellow, who was to become Koki's favourite.





Summer sped by and pre-monsoon showers swept across the hills. The grass on the hillside was transformed from a dull brown to a lush green. Tiny streams became torrents fo rushing water. The insect world came to life. Cicadas sang in the oak trees. A slender praying mantis balanced on the windowsill. At night, months flew into the front room, and sometimes the months flew into the front room, and sometimes the months flew into the goldfish tank and drowned.
\par
Beetles kept blundering in through the open window. They seemed to be attracted by the bulb above the goldfish bowl. Every now in through the open window. They seemed to be attracted by the bulb above the goldfish bowl. Every now and then, one of them would fall into the water, and Koki would have to fish it out. First, an elegant bamboo beetle. Then, next evening, a clumsy bamboo beetle, followed by a stag, floating, struggling, sinking, until Koki came to rescue.
\par
Naturally, the gold fish were alarmed. FLOP! SPLASH! Every now and then, one of those clumsy beetles would land in the tank and be unable to get out. The goldfish were used to stillness, silence and calm water. The beetles were a terrible nuisance, and finally, Koki had to keep the window closed towards evening, when the lights came on.
\par
The goldfish looked grateful.
\par
Koki put her face against the glass of the tank, and her favourite pink goldfish came up to the glass and for a time, they were nose to nose. It might have been trying to say a something. Who knows?
\par
Then came a dog that didn't bark.
\par
Determined that Granny should keep a dog, her sister had sent her a strange-looking hound, long and thin are very hairy, with a note saying it didn't bark but that it would keep the jackals away.
\par
The dog stayed with them for three days without barking. When it spotted squirrels, it whined a little and tried to chase them. And when it saw the cat, it growled and tried to look menaching. But it didn't bark-not even at the postman.
\par
Then, early one right, as the full moon came over the mountain, the dog raised his head and let out the most dreadful howl.
\par
And this set off all the dogs, jackals, wolves and foxes in the forest, and soon they were all howling: some howled at the moon, some howled at each other and some howled because it seemed like the right thing to do.
\par
This went on for about ten minutes and then it stopped just as they suddenly as it started. The moon disappeared behind a cloud and fell silent. All the animals fell asleep. The goldfish didn't seen to mind at all.
\par
Granny, being kind-hearted, didn't send the dog away.
\par
'If it's only the moon that brothers him,' she said, 'we'll just keep the curtains drawn whenever the moon comes up!'
\par
It was that simple.
\par
And the dog would be a good compainon to Granny in the days and months to come.
\par
It was time for Koki to return to her home in the city.
\par
She couldn't take the goldfish with her-not even her little pink favourite one. But she knew she would be back before long, prehaps during the winter break.
\par
She say goodbye to the dog, she said goodbye to the cat, she said goodbye to the birds in the oak trees and she gave Granny for a long, lingering, loving hug.
\par
Then she put her face to the glass of the goldfish tank and said goodbye to all her pretty, silent friend. They didn't say anything. Goldfish don't talk. But the little pink goldfish, Koki's favourite, glided up to the edge of the tank and pressed its face against the glass, right up against Koki's gentle face. Only the glass separated them. Was it trying to say something? Was it saying goodbye? Who knows?
